\section{Rational tensors and the coloring obstruction}\label{sec:tensors}

We first express the matrix condition of Theorem~\ref{thm:classification}
as a condition on affine functions.  This formulation makes the obstruction
to arbitrary colorings explicit and supplies the input for the construction
in the next section.  Throughout, $s\geq2$, $d\geq1$, and
$A=\{a_1,\ldots,a_s\}$ affinely spans $\R^d$; the field $F$, ring $B$,
columns $p_i$, and multiplication map $m_F$ are as in the theorem.

Put $W=\R\times\R^d$, regarded as a vector space over $\Q$.
An element $w=(v,u)$ specifies the affine function $a\mapsto v+a\cdot u$.
For each label define the $\Q$-linear evaluation map
\[
 e_i:W\longrightarrow\R,\qquad e_i(v,u)=v+a_i\cdot u.
\]
Let $\mathcal T(W)$ be the subspace of $W\otimes_{\Q}W$ fixed by the
flip $w\otimes w'\mapsto w'\otimes w$.  Thus our symmetric tensors are
flip-invariant tensors.  The map
\[
 (v,u)\otimes(v',u')\longmapsto u(u')^{\mathsf T}
\]
is $\Q$-linear on the tensor product and restricts to a map
$G:\mathcal T(W)\to\Sym_d(\R)$, where $\Sym_d(\R)$ denotes the real
symmetric matrices.  In particular, $G(w\otimes w)=uu^{\mathsf T}$.
The tensor condition is the existence of $T\in\mathcal T(W)$ such that
\begin{equation}\label{alg:tensor}
 (e_i\otimes e_i)(T)=0\quad(1\leq i\leq s),
 \qquad G(T)=I_d.
\end{equation}
The first identities hold in $\R\otimes_{\Q}\R$.
They are stronger than their images under real multiplication.

\begin{proposition}\label{alg:equivalence}
The matrix condition in Theorem~\ref{thm:classification} is equivalent
to~\eqref{alg:tensor}.
\end{proposition}

\begin{proof}
Write $\mathscr B=\R\otimes_{\Q}\R$, let
$\tau(x\otimes y)=y\otimes x$, and let
$m_{\R}:\mathscr B\to\R$ be multiplication.
The coordinate decomposition of $W$ identifies $W\otimes_{\Q}W$ with
arrays $Q=(Q_{\alpha\beta})_{0\leq\alpha,\beta\leq d}$ over $\mathscr B$.
Under this identification, evaluation is
\[
 E_i(Q)=\sum_{\alpha,\beta=0}^d
 (p_{i,\alpha}\otimes1)Q_{\alpha\beta}
 (1\otimes p_{i,\beta}),
\]
and the gradient matrix is the spatial block
$(m_{\R}(Q_{\alpha\beta}))_{1\leq\alpha,\beta\leq d}$.
Tensor flip becomes
\[
 Q^*_{\alpha\beta}=\tau(Q_{\beta\alpha}).
\]
Directly from the evaluation formula,
$E_i(Q^*)=\tau(E_i(Q))$, and multiplication of $Q^*$ transposes the
multiplied array.  Consequently $\tfrac12(Q+Q^*)$ preserves every zero
evaluation and any multiplied spatial block equal to $I_d$.
This proves the implication from a matrix over $B$ to a symmetric tensor,
using the natural inclusion $B\subset\mathscr B$.

For the converse, we show that extending the coefficient ring from $B$
to $\mathscr B$ introduces no new solutions.  Choose an $F$-basis
$(r_\lambda)_{\lambda\in\Lambda}$ of $\R$.  Distributing tensor products
over direct sums gives the free $B$-module decomposition
\begin{equation}\label{alg:free-expansion}
 \mathscr B=
 \bigoplus_{\lambda,\mu\in\Lambda}
 B(r_\lambda\otimes r_\mu).
\end{equation}
Here a coefficient $x\otimes y\in B$ in the indicated summand represents
$(xr_\lambda)\otimes(yr_\mu)$.
If an array $Q$ over $\mathscr B$ solves the evaluation equations and has
multiplied spatial block $I_d$, its entries have a common finite expansion
\[
 Q=\sum_{(\lambda,\mu)\in J}
 (r_\lambda\otimes r_\mu)P^{\lambda\mu},
 \qquad P^{\lambda\mu}\in\operatorname{Mat}_{d+1}(B).
\]
All coefficients in the maps $E_i$ belong to $B$.
Thus~\eqref{alg:free-expansion} implies
$E_i(P^{\lambda\mu})=0$ for each pair and each label.
Let $C^{\lambda\mu}\in\operatorname{Mat}_d(F)$ be the multiplied spatial
block of $P^{\lambda\mu}$.  Applying multiplication to the expansion gives
\[
 I_d=\sum_{(\lambda,\mu)\in J}
 r_\lambda r_\mu C^{\lambda\mu}.
\]
This finite linear system has coefficients and right-hand side in $F$.
Since it is consistent over $\R$, row reduction over $F$ shows that it
is consistent over $F$.  Choose $c_{\lambda\mu}\in F$ with
$I_d=\sum c_{\lambda\mu}C^{\lambda\mu}$.  Then
\[
 P=\sum_{(\lambda,\mu)\in J}
 (c_{\lambda\mu}\otimes1)P^{\lambda\mu}
\]
is a solution over $B$.  The individual matrices $C^{\lambda\mu}$ need
not be symmetric; the descent took place in the full space
$\operatorname{Mat}_d(F)$.  If a symmetric tensor is wanted after descent,
the symmetrization already proved applies to $P$.
\end{proof}

\begin{remark}\label{alg:noetherian}
The finite coordinate presentation explains the algebraic nature of the
criterion.  Let $F_0=\Q[(a_i)_j]$, the subring of $\R$ generated by all
coordinates.  It is a finitely generated domain and $F=\operatorname{Frac}(F_0)$.
The ring $B$ is the localization of $F\otimes_{\Q}F_0$ at the elements
$1\otimes f$ for $0\ne f\in F_0$.  The latter ring is a finitely generated
$F$-algebra, so $B$ is Noetherian.
The kernel of the finite evaluation system therefore has finitely many
$B$-module generators.  Its multiplied spatial image is the $F$-span of
their multiplied spatial blocks, because every scalar $c\in F$ lifts to
$c\otimes1$.  This description uses exact coordinate-field relations;
it makes no claim to recover those relations from numerical real inputs.
\end{remark}

\begin{remark}\label{alg:congruence}
The tensor condition is directly invariant under congruence.
If $a_i'=Qa_i+t$ with $Q\in O(d)$, the invertible $\Q$-linear map
\[
 H:W\longrightarrow W,\qquad
 H(v,u)=(v-t\cdot Qu,Qu)
\]
satisfies $e_i'\circ H=e_i$ and
$G((H\otimes H)T)=QG(T)Q^{\mathsf T}$.
It therefore carries certificates for $A$ to certificates for $A'$;
the inverse congruence gives the converse.  Reordering labels merely
reorders the equations.  Any two congruent representatives in their
affine dimension are related by such a map, so
Proposition~\ref{alg:equivalence} also proves representative independence
for the matrix criterion, even when their coordinate fields differ.
\end{remark}

We next prove the necessary direction of the classification.
Failure of~\eqref{alg:tensor} gives an algebraic functional, and its values
provide a fixed finite coloring that works in every ambient dimension.
The invariant-and-coloring strategy goes back to Erd\H{o}s et al.
\cite[Theorem~13 and Lemma~15]{egm-rss}; here rational tensor separation
supplies the invariant, and the modular coloring is proved directly.
Their real-variable coloring lemma extends Rado's work on partition
regularity \cite{rado-note}.

\begin{proposition}\label{alg:necessity}
If~\eqref{alg:tensor} has no solution, then there is an integer $r_0\geq2$
such that every $\R^D$ admits an $r_0$-coloring with no monochromatic
congruent copy of $A$.  In particular, every Euclidean Ramsey set
satisfies the matrix condition of Theorem~\ref{thm:classification}.
\end{proposition}

\begin{proof}
Regard
\[
 \mathcal Y=(\R\otimes_{\Q}\R)^s\oplus\Sym_d(\R)
\]
as a $\Q$-vector space and define the $\Q$-linear map
\[
 \Phi:\mathcal T(W)\longrightarrow\mathcal Y,
 \qquad
 \Phi(T)=\bigl((e_i\otimes e_i)(T)\bigr)_{i=1}^s\oplus G(T).
\]
By assumption $y_0=(0,\ldots,0,I_d)$ is outside its image.
In the quotient vector space $\mathcal Y/\operatorname{im}\Phi$, extend
the nonzero class of $y_0$ to a basis over $\Q$.
Sending that basis element to $1$ and the others to $0$ gives a
$\Q$-linear functional $\ell:\mathcal Y\to\R$ such that
\[
 \ell\circ\Phi=0,\qquad \ell(y_0)=1.
\]
Write its components as $\ell_i:\R\otimes_{\Q}\R\to\R$ and
$L:\Sym_d(\R)\to\R$, and put $h_i(x)=\ell_i(x\otimes x)$.
Applying $\ell$ to $\Phi(w\otimes w)$ gives
\begin{equation}\label{alg:separated}
 \sum_{i=1}^s h_i(v+a_i\cdot u)+L(uu^{\mathsf T})=0,
 \qquad L(I_d)=1.
\end{equation}
Taking $u=0$ shows that
\begin{equation}\label{alg:constant-cancellation}
 \sum_{i=1}^s h_i(x)=0\qquad(x\in\R).
\end{equation}
These are purely algebraic functionals; no continuity or measurability is
asserted or needed.

Fix $D$ and define
$H_i(z)=\sum_{j=1}^D h_i(z_j)$ for $z\in\R^D$.
We claim that every ordered congruent copy $(b_1,\ldots,b_s)$ satisfies
\begin{equation}\label{alg:copy-obstruction}
 \sum_{i=1}^s H_i(b_i)=-1.
\end{equation}
Indeed, equality of the pairwise distances gives equality of the Gram
matrices of $(a_i-a_1)_i$ and $(b_i-b_1)_i$ by polarization.
Thus sending each $a_i-a_1$ to $b_i-b_1$ defines an inner-product-preserving
linear map on their span: any relation among the first differences has
squared norm zero for the corresponding combination of the second.
Since $A$ affinely spans $\R^d$, there are a linear isometry
$Q:\R^d\to\R^D$ and $t\in\R^D$ such that $b_i=t+Qa_i$.
Writing $u_j^{\mathsf T}$ for the $j$th row of $Q$, we have
\[
 (b_i)_j=e_i(t_j,u_j),\qquad
 \sum_{j=1}^D u_ju_j^{\mathsf T}=Q^{\mathsf T}Q=I_d.
\]
Sum~\eqref{alg:separated} over $j$.
The finite additivity of $L$ proves~\eqref{alg:copy-obstruction}.
Also~\eqref{alg:constant-cancellation} gives
$\sum_iH_i(z)=0$ at every single point $z$.

Choose $K=2s+1$ and partition $[0,2)$ into $K$ half-open intervals
of equal length $2/K<1/s$.  Color a point $z$ by the $s$-tuple of interval
labels of the representatives of $H_i(z)$ modulo $2$.
This uses at most $r_0=K^s$ colors, independently of $D$.
If $b_1,\ldots,b_s$ were monochromatic, there would be integers $n_i$
and real numbers $\varepsilon_i$, with $|\varepsilon_i|<2/K$, such that
\[
 H_i(b_i)-H_i(b_1)=2n_i+\varepsilon_i.
\]
Summation yields
\[
 -1=2\sum_i n_i+\sum_i\varepsilon_i,
 \qquad \left|\sum_i\varepsilon_i\right|<1.
\]
The first identity makes the error sum an odd integer, contradicting the
second.  This coloring avoids $A$ in every dimension and completes the
proof.
\end{proof}

The remaining direction begins with a tensor satisfying
\eqref{alg:tensor}.  Its evaluations vanish before multiplication, and
that exact cancellation will produce equal coordinate histograms for
two embeddings at different scales.
